%------------------------------ Ejercicio 5.30 -------------------------------
\setcounter{neje}{29}
\begin{nejer}
Sea $X \coloneq \mathbb{R}$, denote por $\mathcal{B} \subseteq 2^X$ a la $\sigma$-álgebra de Borel, y sea $\mu \colon \mathcal{B} \to [0, \infty]$ la restricción de la medida de Lebesgue a $\mathcal{B}$. Sea $\lambda \colon \mathcal{B} \to [0, \infty]$ una medida. Pruebe lo siguiente.
\begin{enumerate}[label = (\roman*), leftmargin=1.8em]
    \item Si $B \en \mathcal{B}$ y $0 < c < \mu(B)$, entonces existe un conjunto de Borel $A \subseteq B$ tal que $\mu(A) = c$. \emph{Sugerencia.} Muestre que la función $f(t) \coloneq \mu(B \cap [-t, t])$ es continua.
    \item Si existe una constante $0 < c < \infty$ tal que
    $$
    \mu(B) = c \implies \lambda(B) = c
    $$
    para todo $B \en \mathcal{B}$, entonces $\lambda \ll \mu$.
\end{enumerate}
\end{nejer}
\begin{sol}
\begin{enumerate}[label=(\roman*), leftmargin=1.8em]
%--------------------------------------------------
\item % (i)
\emph{Veamos que $f\colon \mathbb{R}_{\ge0} \to \mathbb{R}$ tal que $f(t) \coloneq \mu(B\cap[-t,t])$ es 2-Lipschitz}:\\
Sean $t_1,t_2\en\mathbb{R}_{\ge0}$ arbitrarios, y supongamos sin perder generalidad que $t_1 \leq t_2$. Luego, gracias a que $B \cap [-t_1,t_1] \subseteq B \cap [-t_2,t_2]$ y $\mu(B \cap [-t_2,t_2]) < \infty$, la aditividad garantiza:
\[
\abs{ 
    \mu\bigl( B \cap [-t_2,t_2] \bigr) - \mu\bigl( B \cap [-t_1,t_1] \bigr)
}
\;=\; \mu\bigl((B \cap [-t_2,t_2]) \setminus(B \cap [-t_1,t_1])\bigr)
\]
pero es inmediato que $(B \cap [-t_2,t_2]) \setminus(B \cap [-t_1,t_1]) = B \cap [-t_2,t_2]\setminus[-t_1,t_1]$. Así que:
\begin{align*}
\abs{ 
    \mu\bigl( B \cap [-t_2,t_2] \bigr) - \mu\bigl( B \cap [-t_1,t_1] \bigr)
}
&\;=\; \mu\bigl(B \cap [-t_2,t_2] \setminus [-t_1,t_1]\bigr)\\
&\;\leq\; \mu\bigl([-t_2,t_2] \setminus [-t_1,t_1]\bigr)\\
&\;=\; \mu\bigl([-t_2,-t_1) \cup (t_1,t_2]\bigr)\;=\; 2(t_2-t_1)
\end{align*}
En otras palabras, $\abs{f(t_2) - f(t_1)} \leq 2\abs{t_2-t_1}$, donde $t_1,t_2\en\mathbb{R}_{\ge0}$ son arbitrarios.\par
Ahora bien, como $\set{B\cap[-n,n]}_{n\en\mathbb{N}}$ es monótona con unión $B$, tenemos
\[
\lim_{n\to\infty}f(n) 
= \lim_{n\to\infty}\mu(B\cap[-n,n])
=\mu(B)
\]
Por consiguiente, si $0<c<\mu(B)$, forzosamente existe $n_0\en\mathbb{N}$ tal que $c < f(n_0) \le \mu(B)$; sin embargo, como
$f$ es Lipschitz, es (uniformemente) continua, así que $f(0) = 0 < c < f(n_0)$ garantiza (por teo. del valor intermedio) que existe $t_0\en(0,n_0)$ tal que $f(t_0) = c$. En otras palabras, $A \coloneq B \cap [-t_0,t_0]$ satisface $\mu(A) = c$.
%--------------------------------------------------
\item % (ii)
Supongamos que $c\en\mathbb{R}^+$ es tal que "$\mu(B) = c \implies \lambda(B)=c$".
Luego, sea $A\en\mathcal{A}$ arbitrario con $\mu(A)=0$, buscamos probar que $\lambda(A) = 0$. En efecto, por aditividad 
\[
\mu(A^c) \;=\; \mu(A) + \mu(A^c) \;=\; \mu(X) \;=\; \infty
\] 
por lo que $0<c<\mu(A^c)$, y así, el inciso anterior garantiza que existe $C\subseteq A^c$ medible con $\mu(C) = c$, de donde también $\lambda(C) = c$. Luego, como $A \cap C = \varnothing$, 
\[
\mu(A\cup C) \;=\; \mu(A) + \mu(C) \;=\; \mu(C) \;=\; c
\]
de donde se sigue también $\lambda(A \cup C) = c$. Solo resta aplicar aditividad de $\lambda$:
\[
c 
= \lambda(A \cup C) = \lambda(A) + \lambda(C) = \lambda(A) + c
\qquad\text{lo cual implica}\qquad
\lambda(A) = 0
\]
Como $A\en\mathcal{A}$ es arbitrario, esto prueba que $\lambda \ll \mu$.
\end{enumerate}
\end{sol}